% homework/hw07x/hw07x.tex
% Year 7 - Extra Homework: split, multiply, add.
%
% Goes out alongside Homework 7. Same framing rule as hw03x through hw06x:
% nothing on the sheet says easy, basic, extra help, catch-up or revision, and
% no question is marked as being for anybody in particular. The cover states a
% fact, and it is a true one:
%
%   two-digit multiplication and expanding a bracket are the same picture.
%
% THAT IS THE WHOLE DESIGN. 4 x 16 is 4(10 + 6) = 40 + 24, which is the deck's
% own starter for 2.4, and 4(x + 3) = 4x + 12 is the same grid with one box
% unknown. So this sheet gives the two-digit multiplication practice the class
% needs AND lands them on 2.4's method, rather than treating the arithmetic and
% the algebra as two subjects.
%
% hw05x already taught the column method -- two rows, then add -- and long
% division. This sheet does not repeat either. It teaches the GRID, which the
% column method hides and the algebra needs, and its division is aimed at one
% thing: the reverse of a multiplication, which is what undoes a bracket.
%
% F20 is the sheet in miniature and the last question on purpose. Expand
% 4(x + 3), then put x = 10 into both sides: 4 x 13 = 52 and 4 x 10 + 12 = 52.
% The arithmetic they practised in Part 1 proves the algebra in Part 4, using
% the same grid both times.
%
% House style is shared: see homework/hw-style.tex. Method: the new-homework
% skill. Source is pure ASCII on purpose.
%
% Build:  pwsh .claude/skills/new-homework/scripts/build-hw.ps1 hw07x

\documentclass[12pt,a4paper]{article}

% -- Answer key switch -------------------------------------------------------
\newif\ifkey
\ifdefined\TEACHER\keytrue\else
  \keyfalse        % \keyfalse = student copy   |   \keytrue = teacher copy
\fi

% -- House style -------------------------------------------------------------
% homework/hw-style.tex
% Shared house style for every Year 7 homework packet. \input this from the
% packet file, right after \documentclass and the answer-key switch.
%
% Do not fork this file per packet. If a packet needs a different look, the
% question is whether every packet should have it.
%
% The choices here are deliberate and were arrived at by printing and looking:
%   * No solid dark banners. Every header is a light tint with a coloured spine
%     on the left, so colour reads clearly without flooding a school printer.
%   * No marks and no mark totals. These are practice, not exam papers.
%   * Lato at 12pt with open leading, plus mathastext so inline maths is Lato
%     too. Without mathastext every "-6 + 4" is Computer Modern serif and the
%     sheet looks half-typeset.
%   * Writing lines are spaced for Year 7 handwriting, not adult margin notes.
%
% Provides: \wlines \blank \tickbox \sep \qmk-free marks (none) \hwsection
% \qhead \expr, the workspace box, and the remember / wordhelp / activity /
% yourturn coloured boxes. Palette matches the lesson decks exactly.

% -- Packages ----------------------------------------------------------------
\usepackage[T1]{fontenc}
\usepackage[utf8]{inputenc}
\usepackage[default]{lato}
\usepackage[a4paper,top=1.7cm,bottom=1.7cm,left=2.0cm,right=2.0cm,headheight=15pt]{geometry}
\usepackage{amsmath,amssymb}
\usepackage{xcolor}
\usepackage[most]{tcolorbox}
\usepackage{enumitem}
\usepackage{tikz}
\usetikzlibrary{arrows.meta}
\usepackage{array}
\usepackage{tabularx}
\usepackage{fancyhdr}
\usepackage{multicol}
\usepackage{needspace}
\usepackage{graphicx}

% Make maths use Lato too. Without this, every "-6 + 4" on the page is set in
% Computer Modern serif and the sheet looks half-typeset. Must come last.
\usepackage[italic]{mathastext}

\linespread{1.06}\selectfont

% -- The Learner's Book palette (same hex values as the lesson decks) ---------
\definecolor{cbteal}{HTML}{0087A8}
\definecolor{cbpurple}{HTML}{5C2483}
\definecolor{cborange}{HTML}{C25E12}
\definecolor{cbgreen}{HTML}{4A8B23}
\definecolor{cbred}{HTML}{C8102E}
\definecolor{cbblue}{HTML}{1A5FA8}
\definecolor{rulegray}{HTML}{9AA5AD}

% -- Answer lines and blanks -------------------------------------------------
% \wlines{n} draws n ruled writing lines, spaced wide enough for Year 7
% handwriting rather than for an adult with a fine pen.
\makeatletter
\newcommand{\wlines}[1]{%
  \par\vspace{0.75em}%
  \@tempcnta=0\relax
  \@whilenum\@tempcnta<#1\do{%
    \hbox to \linewidth{\color{rulegray}\leaders\hrule height 0.5pt\hfill}%
    \vspace{1.5em}%
    \advance\@tempcnta by 1\relax}%
  \par\vspace{-0.8em}%
}
\makeatother

% \blank[width] is an inline gap to write a short answer in.
\newcommand{\blank}[1][2.6cm]{\underline{\hspace{#1}}}

% A boxed space for showing working.
\newtcolorbox{workspace}[1][3.4cm]{%
  colback=white, colframe=rulegray, boxrule=0.5pt, arc=5pt,
  height=#1, left=6pt, right=6pt, top=4pt, bottom=4pt,
  before skip=9pt, after skip=11pt}

% An empty tick box.
\newcommand{\tickbox}{\fbox{\rule{0pt}{1.15em}\rule{1.15em}{0pt}}}

% A middot separator.
\newcommand{\sep}{\,\textperiodcentered\,}

% -- Section and question headings -------------------------------------------
% Light tint plus a fat coloured spine on the left. No solid dark banner: the
% colour reads clearly but barely costs the printer anything.
% needspace stops a heading stranding alone at the foot of a page. Kept modest:
% too large and it pushes headings over, leaving a worse hole than it prevents.
\newcommand{\hwsection}[2][cbteal]{%
  \needspace{8\baselineskip}%
  \par\vspace{13pt}%
  \noindent
  \begin{tcolorbox}[enhanced, nobeforeafter, width=\textwidth,
      colback=#1!7, colframe=#1!7, arc=5pt,
      borderline west={5pt}{0pt}{#1},
      left=14pt, right=10pt, top=7pt, bottom=7pt]
    {\large\bfseries\color{#1}#2}
  \end{tcolorbox}%
  \par\vspace{9pt}}

% \qhead[colour]{A2}{Moving along the line} -- a question heading.
\newcommand{\qhead}[3][cbteal]{%
  \needspace{6\baselineskip}%
  \par\vspace{11pt}%
  \noindent{\large\bfseries\color{#1}#2}\quad{\large\bfseries #3}\par\vspace{4pt}}

% -- Coloured boxes ----------------------------------------------------------
% Shared look: rounded, light fill, coloured rule, coloured title text on a
% pale title strip rather than a dark one.
\tcbset{hwbox/.style={
  enhanced, breakable, arc=6pt, boxrule=1pt,
  left=11pt, right=11pt, top=8pt, bottom=8pt,
  fonttitle=\bfseries, attach boxed title to top left={xshift=12pt, yshift=-8pt},
  boxed title style={arc=4pt, boxrule=1pt, left=7pt, right=7pt, top=2pt, bottom=2pt},
  top=13pt}}

% Orange: things to remember. Matches the deck's "Write This Down".
\newtcolorbox{remember}[1][Remember from the lesson]{hwbox,
  colback=cborange!6, colframe=cborange, coltitle=cborange,
  colbacktitle=white, title={#1}}

% Blue: English help. The barrier in this class is the wording, not the maths.
\newtcolorbox{wordhelp}[1][Word help]{hwbox,
  colback=cbblue!6, colframe=cbblue, coltitle=cbblue,
  colbacktitle=white, title={#1}}

% Purple: an activity or instruction.
\newtcolorbox{activity}[1][How to do this homework]{hwbox,
  colback=cbpurple!6, colframe=cbpurple, coltitle=cbpurple,
  colbacktitle=white, title={#1}}

% Green: your turn / a friendly nudge.
\newtcolorbox{yourturn}[1][Your turn]{hwbox,
  colback=cbgreen!7, colframe=cbgreen, coltitle=cbgreen!75!black,
  colbacktitle=white, title={#1}}

% -- Shorthands --------------------------------------------------------------
\newcommand{\dC}{\,^{\circ}\mathrm{C}}          % use inside maths: $-8\dC$
\newcommand{\key}[1]{{\color{cborange}\textbf{#1}}}
\newcommand{\eng}[1]{{\color{cbblue}\textbf{#1}}}

\setlist[enumerate]{itemsep=8pt, topsep=5pt, leftmargin=*}
\setlist[itemize]{itemsep=4pt, topsep=4pt, leftmargin=*}
\renewcommand{\arraystretch}{1.45}

% -- An expression the student is asked to work from -------------------------
\newcommand{\expr}[1]{%
  \tcbox[on line, colback=cbgreen!10, colframe=cbgreen!55, boxrule=1pt, arc=4pt,
         left=9pt, right=9pt, top=4pt, bottom=4pt]{\large\bfseries $#1$}}

% -- Running head and foot ---------------------------------------------------
% The packet overrides \fancyhead[L] with its own title.
\pagestyle{fancy}
\fancyhf{}
\fancyhead[L]{\small\color{black!55}Year 7\sep Homework}
\fancyhead[R]{\small\color{black!55}Name: \blank[3.4cm]}
\fancyfoot[C]{\small\color{black!55}Page \thepage}
\renewcommand{\headrulewidth}{0pt}
\renewcommand{\headrule}{\vspace{2pt}\hbox to\headwidth{\color{cbteal!45}\leaders\hrule height 1.2pt\hfill}}



% -- Per-packet running head -------------------------------------------------
\fancyhead[L]{\small\color{black!55}Year 7\sep Extra Homework}

% -- A group of writing lines that will not be split -------------------------
% See hw06.tex: \wlines is a run of separate hboxes and TeX will break in the
% middle of one, stranding a single rule at the top of the next page.
\newcommand{\qlines}[1]{\par\needspace{\numexpr#1+1\relax\baselineskip}\wlines{#1}}

% -- A box to write one number in --------------------------------------------
% The workbook's own scaffold: a part-written line with an empty square where
% the number goes. Filling the squares IS the method being taught, so it has to
% look like something you write in, not like a blank.
\newcommand{\nbox}{\ \fbox{\rule{0pt}{1.05em}\rule{1.0em}{0pt}}\ }

% -- An empty two-by-two grid to multiply in ---------------------------------
% #1 #2 are the column headings (the number split into tens and ones), #3 #4
% the row headings. Every practice grid on the sheet is one of these, so the
% picture never changes between a number question and a letter question -- which
% is the entire argument of the sheet.
\newcommand{\gridtwo}[4]{%
  \renewcommand{\arraystretch}{1.9}%
  \begin{tabular}{|c|c|c|}
    \hline
    $\times$ & $#1$ & $#2$ \\ \hline
    $#3$     &      &      \\ \hline
    $#4$     &      &      \\ \hline
  \end{tabular}%
  \renewcommand{\arraystretch}{1.45}}

% -- An empty one-row grid ---------------------------------------------------
\newcommand{\gridone}[3]{%
  \renewcommand{\arraystretch}{1.9}%
  \begin{tabular}{|c|c|c|}
    \hline
    $\times$ & $#1$ & $#2$ \\ \hline
    $#3$     &      &      \\ \hline
  \end{tabular}%
  \renewcommand{\arraystretch}{1.45}}

\begin{document}

% ===========================================================================
% COVER BLOCK
% ===========================================================================
\thispagestyle{fancy}

\begin{tcolorbox}[enhanced, colback=cbgreen!8, colframe=cbgreen!8, arc=7pt,
                  borderline west={6pt}{0pt}{cbgreen},
                  left=18pt, right=14pt, top=11pt, bottom=11pt]
  {\color{cbgreen!75!black}\bfseries YEAR 7\sep UNIT 2\sep EXTRA HOMEWORK}\\[6pt]
  {\color{cbgreen!70!black}\Huge\bfseries Split, Multiply, Add}\\[10pt]
  {\color{black!70}Two-digit multiplication\sep dividing\sep and the same grid with a letter in it}
\end{tcolorbox}

\vspace{6pt}
\noindent
\begin{tabularx}{\textwidth}{@{}X X@{}}
  \textbf{Name:} \blank[5.6cm] & \textbf{Class:} \blank[5.6cm] \\
  \textbf{Date given:} \blank[4.6cm] &
  \textbf{Date due:} {\color{cbred}\bfseries Monday 21 September 2026} \\
\end{tabularx}

\vspace{4pt}
\begin{activity}[Why this sheet exists]
  Nobody multiplies by 16. They \key{split it up}: $4 \times 10$ is 40,
  $4 \times 6$ is 24, and $40 + 24 = 64$.

  \medskip
  That is three steps --- \key{split, multiply, add} --- and it is also, exactly,
  what expanding a bracket is. $4 \times 16$ and $4(x + 3)$ are drawn in the
  \key{same grid}. Part 1 of this sheet practises the grid on numbers. Part 4
  puts a letter in one box. Nothing else changes.

  \medskip
  \eng{No calculator.} Write in every box.
\end{activity}

% ===========================================================================
% PART 1 -- SPLIT IT UP
% ===========================================================================
\hwsection[cbgreen]{Part 1 \textperiodcentered\ Split it up}

\qhead[cbgreen]{F1}{Tens and ones}

\noindent Split each number. The first one is done for you.

\vspace{4pt}
\begin{multicols}{3}
\begin{enumerate}[label=(\alph*), itemsep=10pt]
  \item $47 = 40 + 7$
  \item $23 = \blank[2.2cm]$
  \item $65 = \blank[2.2cm]$
  \item $18 = \blank[2.2cm]$
  \item $52 = \blank[2.2cm]$
  \item $96 = \blank[2.2cm]$
  \item $34 = \blank[2.2cm]$
  \item $70 = \blank[2.2cm]$
  \item $89 = \blank[2.2cm]$
\end{enumerate}
\end{multicols}

% -- F2 ---------------------------------------------------------------------
\qhead[cbgreen]{F2}{Times a ten}

\noindent This is the box that people get wrong, so do a page of them first.

\begin{remember}[Multiply, then put the zero back on]
  \hspace*{1.2em}$4 \times 30$: \; $4 \times 3 = 12$, \; so \; $4 \times 30 =
  \mathbf{120}$.

  \smallskip
  \hspace*{1.2em}$20 \times 30$: \; $2 \times 3 = 6$, \; \key{two} zeros, \; so \;
  $20 \times 30 = \mathbf{600}$.
\end{remember}

\vspace{2pt}
\begin{multicols}{4}
\begin{enumerate}[label=(\alph*), itemsep=10pt]
  \item $4 \times 30 = \blank[1.3cm]$
  \item $6 \times 50 = \blank[1.3cm]$
  \item $7 \times 40 = \blank[1.3cm]$
  \item $9 \times 60 = \blank[1.3cm]$
  \item $8 \times 20 = \blank[1.3cm]$
  \item $3 \times 90 = \blank[1.3cm]$
  \item $20 \times 30 = \blank[1.3cm]$
  \item $40 \times 50 = \blank[1.3cm]$
  \item $60 \times 70 = \blank[1.3cm]$
  \item $30 \times 80 = \blank[1.3cm]$
  \item $50 \times 50 = \blank[1.3cm]$
  \item $20 \times 90 = \blank[1.3cm]$
\end{enumerate}
\end{multicols}

% -- F3 ---------------------------------------------------------------------
\qhead[cbgreen]{F3}{One row: two digits by one digit}

\noindent Fill the two boxes, then add them.

\begin{remember}[The grid, once, in full]
  $4 \times 16$ \; means \; $4 \times (10 + 6)$. \; One box, one multiplication.

  \medskip
  \noindent
  \begin{minipage}[c]{0.40\linewidth}
    \renewcommand{\arraystretch}{1.9}%
    \begin{tabular}{|c|c|c|}
      \hline
      $\times$ & $10$ & $6$ \\ \hline
      $4$      & $40$ & $24$ \\ \hline
    \end{tabular}%
    \renewcommand{\arraystretch}{1.45}
  \end{minipage}%
  \begin{minipage}[c]{0.58\linewidth}
    $40 + 24 = \mathbf{64}$

    \smallskip
    \key{Split} 16 into $10 + 6$. \; \key{Multiply} each box. \; \key{Add} the
    boxes.
  \end{minipage}
\end{remember}

\vspace{4pt}
\noindent
\begin{minipage}[c]{0.30\textwidth}\gridone{20}{3}{7}\end{minipage}%
\begin{minipage}[c]{0.20\textwidth}\textbf{(a)} $7 \times 23$ \\[4pt] $= \blank[1.8cm]$\end{minipage}%
\begin{minipage}[c]{0.30\textwidth}\gridone{40}{8}{6}\end{minipage}%
\begin{minipage}[c]{0.20\textwidth}\textbf{(b)} $6 \times 48$ \\[4pt] $= \blank[1.8cm]$\end{minipage}

\vspace{11pt}
\noindent
\begin{minipage}[c]{0.30\textwidth}\gridone{50}{4}{8}\end{minipage}%
\begin{minipage}[c]{0.20\textwidth}\textbf{(c)} $8 \times 54$ \\[4pt] $= \blank[1.8cm]$\end{minipage}%
\begin{minipage}[c]{0.30\textwidth}\gridone{30}{7}{9}\end{minipage}%
\begin{minipage}[c]{0.20\textwidth}\textbf{(d)} $9 \times 37$ \\[4pt] $= \blank[1.8cm]$\end{minipage}

% -- F4 ---------------------------------------------------------------------
\qhead[cbgreen]{F4}{Now without the printed grid}

\noindent Draw your own grid in the box, then write the answer under it.

\vspace{4pt}
\noindent
% Three bands rather than one table, for the reason given at F6.
\begin{tabularx}{\textwidth}{|X|X|X|X|}
  \hline
  $6 \times 34$ & $8 \times 72$ & $7 \times 85$ & $9 \times 46$ \\
  \rule[-3.0em]{0pt}{3.6em} & & & \\
  \hline
\end{tabularx}

\vspace{5pt}
\noindent
\begin{tabularx}{\textwidth}{|X|X|X|X|}
  \hline
  $5 \times 68$ & $4 \times 93$ & $8 \times 57$ & $7 \times 64$ \\
  \rule[-3.0em]{0pt}{3.6em} & & & \\
  \hline
\end{tabularx}

\vspace{5pt}
\noindent
\begin{tabularx}{\textwidth}{|X|X|X|X|}
  \hline
  $9 \times 78$ & $6 \times 95$ & $7 \times 49$ & $8 \times 86$ \\
  \rule[-3.0em]{0pt}{3.6em} & & & \\
  \hline
\end{tabularx}

% ===========================================================================
% PART 2 -- TWO DIGITS BY TWO DIGITS
% ===========================================================================
\hwsection[cbgreen]{Part 2 \textperiodcentered\ Two digits by two digits}

\qhead[cbgreen]{F5}{Four boxes instead of two}

\noindent Both numbers get split, so the grid has \textbf{four} boxes.

\begin{remember}[$23 \times 14$, once, in full]
  \noindent
  \begin{minipage}[c]{0.40\linewidth}
    \renewcommand{\arraystretch}{1.9}%
    \begin{tabular}{|c|c|c|}
      \hline
      $\times$ & $20$  & $3$ \\ \hline
      $10$     & $200$ & $30$ \\ \hline
      $4$      & $80$  & $12$ \\ \hline
    \end{tabular}%
    \renewcommand{\arraystretch}{1.45}
  \end{minipage}%
  \begin{minipage}[c]{0.58\linewidth}
    $200 + 30 + 80 + 12 = \mathbf{322}$

    \smallskip
    \eng{Four boxes, four multiplications, then add all four.} Three boxes means
    one was missed, and the answer comes out far too small.
  \end{minipage}
\end{remember}

\vspace{4pt}
\noindent
\begin{minipage}[c]{0.30\textwidth}\gridtwo{30}{2}{10}{5}\end{minipage}%
\begin{minipage}[c]{0.20\textwidth}\textbf{(a)} $32 \times 15$ \\[4pt] $= \blank[1.8cm]$\end{minipage}%
\begin{minipage}[c]{0.30\textwidth}\gridtwo{40}{6}{20}{3}\end{minipage}%
\begin{minipage}[c]{0.20\textwidth}\textbf{(b)} $46 \times 23$ \\[4pt] $= \blank[1.8cm]$\end{minipage}

\vspace{11pt}
\noindent
\begin{minipage}[c]{0.30\textwidth}\gridtwo{50}{7}{30}{4}\end{minipage}%
\begin{minipage}[c]{0.20\textwidth}\textbf{(c)} $57 \times 34$ \\[4pt] $= \blank[1.8cm]$\end{minipage}%
\begin{minipage}[c]{0.30\textwidth}\gridtwo{60}{8}{40}{2}\end{minipage}%
\begin{minipage}[c]{0.20\textwidth}\textbf{(d)} $68 \times 42$ \\[4pt] $= \blank[1.8cm]$\end{minipage}

\vspace{11pt}
\noindent
\begin{minipage}[c]{0.30\textwidth}\gridtwo{70}{4}{20}{6}\end{minipage}%
\begin{minipage}[c]{0.20\textwidth}\textbf{(e)} $74 \times 26$ \\[4pt] $= \blank[1.8cm]$\end{minipage}%
\begin{minipage}[c]{0.30\textwidth}\gridtwo{80}{5}{30}{7}\end{minipage}%
\begin{minipage}[c]{0.20\textwidth}\textbf{(f)} $85 \times 37$ \\[4pt] $= \blank[1.8cm]$\end{minipage}

% -- F6 ---------------------------------------------------------------------
\qhead[cbgreen]{F6}{Draw the grid yourself}

\noindent Twelve of them. Draw a four-box grid in each square.

% Four separate tables of three, not one table of twelve. A tabularx cannot be
% broken across a page, so as a single block this walked whole to the next page
% and left two-fifths of the previous one white. As four bands, TeX can break
% between any two of them, and the student sees no difference.
\vspace{4pt}
\noindent
\begin{tabularx}{\textwidth}{|X|X|X|}
  \hline
  $24 \times 13$ & $35 \times 21$ & $47 \times 16$ \\
  \rule[-3.6em]{0pt}{4.2em} & & \\
  \hline
\end{tabularx}

\vspace{5pt}
\noindent
\begin{tabularx}{\textwidth}{|X|X|X|}
  \hline
  $52 \times 34$ & $63 \times 27$ & $78 \times 45$ \\
  \rule[-3.6em]{0pt}{4.2em} & & \\
  \hline
\end{tabularx}

\vspace{5pt}
\noindent
\begin{tabularx}{\textwidth}{|X|X|X|}
  \hline
  $86 \times 53$ & $94 \times 38$ & $75 \times 64$ \\
  \rule[-3.6em]{0pt}{4.2em} & & \\
  \hline
\end{tabularx}

\vspace{5pt}
\noindent
\begin{tabularx}{\textwidth}{|X|X|X|}
  \hline
  $29 \times 37$ & $56 \times 48$ & $83 \times 26$ \\
  \rule[-3.6em]{0pt}{4.2em} & & \\
  \hline
\end{tabularx}

% -- F7 ---------------------------------------------------------------------
\qhead[cbgreen]{F7}{Use it}

\begin{wordhelp}[Word help --- the words that mean multiply]
  \begin{tabularx}{\linewidth}{@{}l@{\hspace{16pt}}X@{}}
    \eng{each} \sep \eng{every} \sep \eng{per} & one of them costs this much, so all of them cost this much times the number \\
    \eng{altogether} \sep \eng{in total} & the whole lot added up \\
    \eng{rows of} \sep \eng{boxes of} & this many groups, each the same size \\
  \end{tabularx}
\end{wordhelp}

\begin{enumerate}[label=\textbf{\arabic*.}, itemsep=10pt]
  \item Mr Bowen buys \textbf{24} boxes of pencils. Every box holds \textbf{36}
        pencils. How many pencils altogether?
        \begin{workspace}[2.2cm]\end{workspace}
  \item A classroom has \textbf{18} rows of chairs with \textbf{15} chairs in each
        row. How many chairs are there?
        \begin{workspace}[2.2cm]\end{workspace}
  \item A bus has \textbf{47} seats. A school books \textbf{23} buses for a trip.
        How many seats are there in total?
        \begin{workspace}[2.2cm]\end{workspace}
  \item Mr Bowen keeps \textbf{18} fish tanks. Every tank holds \textbf{25} rubber
        ducks. How many rubber ducks does he own?
        \begin{workspace}[2.2cm]\end{workspace}
\end{enumerate}

% ===========================================================================
% PART 3 -- DIVIDING
% ===========================================================================
\hwsection[cbgreen]{Part 3 \textperiodcentered\ Dividing --- the multiplication backwards}

\qhead[cbgreen]{F8}{The same three numbers, four ways}

\noindent Every multiplication you know is also two divisions. The first one is done
for you.

\begin{wordhelp}[Word help --- the words that mean divide]
  \begin{tabularx}{\linewidth}{@{}l@{\hspace{16pt}}X@{}}
    \eng{share} \sep \eng{split between} & divide it up so everybody gets the same \\
    \eng{goes into}  & ``6 goes into 54 nine times'' means $54 \div 6 = 9$ \\
    \eng{how many each} & the answer to a division \\
    \eng{left over} \sep \eng{remainder} & what will not divide up evenly \\
  \end{tabularx}
\end{wordhelp}

\vspace{2pt}
\noindent
\begin{tabularx}{\textwidth}{|X|X|X|X|}
  \hline
  \textbf{Multiplication} & \textbf{The other one} & \textbf{Division} &
  \textbf{The other one} \\
  \hline
  $6 \times 9 = 54$ \rule[-0.7em]{0pt}{2.2em} & $9 \times 6 = 54$ &
  $54 \div 6 = 9$ & $54 \div 9 = 6$ \\ \hline
  $7 \times 8 = 56$ \rule[-0.7em]{0pt}{2.2em} & & & \\ \hline
  \rule[-0.7em]{0pt}{2.2em} & & $72 \div 8 = \blank[1.0cm]$ & \\ \hline
  \rule[-0.7em]{0pt}{2.2em} & & & $96 \div 12 = \blank[1.0cm]$ \\ \hline
\end{tabularx}

% -- F9 ---------------------------------------------------------------------
\qhead[cbgreen]{F9}{Divide}

\vspace{2pt}
\begin{multicols}{4}
\begin{enumerate}[label=(\alph*), itemsep=10pt]
  \item $84 \div 7 = \blank[1.2cm]$
  \item $96 \div 8 = \blank[1.2cm]$
  \item $72 \div 6 = \blank[1.2cm]$
  \item $135 \div 9 = \blank[1.2cm]$
  \item $128 \div 4 = \blank[1.2cm]$
  \item $161 \div 7 = \blank[1.2cm]$
  \item $192 \div 8 = \blank[1.2cm]$
  \item $216 \div 6 = \blank[1.2cm]$
\end{enumerate}
\end{multicols}

\vspace{2pt}
\noindent Now four with something \textbf{left over}. Write them like this:
\; $75 \div 8 = 9$ r 3.

\vspace{4pt}
% Two columns, not four. An answer here is "35 r 3", so the blank has to be wide
% enough to write two numbers and a letter in -- and at four columns the blank had
% nowhere to go and dropped onto a line of its own under the division.
\begin{multicols}{2}
\begin{enumerate}[label=(\alph*), start=9, itemsep=11pt]
  \item $58 \div 7 = \blank[2.6cm]$
  \item $100 \div 6 = \blank[2.6cm]$
  \item $85 \div 9 = \blank[2.6cm]$
  \item $143 \div 4 = \blank[2.6cm]$
\end{enumerate}
\end{multicols}

% -- F10 --------------------------------------------------------------------
% The bridge question. Every box here is filled by dividing, and F19 is this
% same question with a bracket around it.
\qhead[cbgreen]{F10}{Find the missing number}

\noindent Every box holds one number. To find it, \key{divide}.

\vspace{4pt}
\begin{multicols}{3}
\begin{enumerate}[label=(\alph*), itemsep=12pt]
  \item $6 \times \nbox = 54$
  \item $\nbox \times 7 = 84$
  \item $9 \times \nbox = 108$
  \item $\nbox \times 8 = 96$
  \item $12 \times \nbox = 60$
  \item $\nbox \times 5 = 115$
\end{enumerate}
\end{multicols}

% -- F11 --------------------------------------------------------------------
\qhead[cbgreen]{F11}{Share it out}

\begin{enumerate}[label=\textbf{\arabic*.}, itemsep=10pt]
  \item \textbf{144} pencils are shared equally between \textbf{12} students. How
        many does each student get?
        \begin{workspace}[2.0cm]\end{workspace}
  \item \textbf{90} students go on a trip. Each bus holds \textbf{16} students.
        How many buses are needed? \; \eng{Careful:} the last bus will not be full,
        and nobody can be left behind.
        \begin{workspace}[2.2cm]\end{workspace}
  \item Mr Bowen shares \textbf{216} stickers equally between \textbf{8} classes.
        How many stickers does each class get?
        \begin{workspace}[2.2cm]\end{workspace}
  \item Mr Bowen buys \textbf{175} rubber ducks and puts the same number into each
        of his \textbf{18} fish tanks. How many go into each tank? How many are
        left over?
        \begin{workspace}[2.4cm]\end{workspace}
\end{enumerate}

% ===========================================================================
% PART 4 -- THE SAME GRID, WITH A LETTER
% ===========================================================================
\hwsection[cbgreen]{Part 4 \textperiodcentered\ The same grid, with a letter in it}

\qhead[cbgreen]{F12}{A letter is a number you do not know yet}

\noindent Collect the like terms. The first one is done for you.

\begin{remember}[Two rules, and that is all of it]
  \key{Only the same letter can be added.} \; $3a + 2a = 5a$, \; but $3a + 2b$
  \key{cannot be added} and is already finished.

  \smallskip
  \key{A letter on its own means one of it.} \; $5x - x = 5x - 1x = \mathbf{4x}$,
  \; {\color{cbred}\bfseries not 5}.
\end{remember}

\vspace{2pt}
\begin{multicols}{3}
\begin{enumerate}[label=(\alph*), itemsep=11pt]
  \item $3a + 2a = 5a$
  \item $4b + 5b = \blank[1.4cm]$
  \item $7c - 2c = \blank[1.4cm]$
  \item $5x - x = \blank[1.4cm]$
  \item $m + m = \blank[1.4cm]$
  \item $6y - y = \blank[1.4cm]$
  \item $2p + 3q = \blank[1.4cm]$
  \item $8k - 3k = \blank[1.4cm]$
  \item $4n + 7 = \blank[1.4cm]$
\end{enumerate}
\end{multicols}

% -- F13 --------------------------------------------------------------------
\qhead[cbgreen]{F13}{One box has a letter in it}

\noindent Exactly the grid from Part 1. One box is a letter, so you cannot work that
box out --- \key{leave it as it is}.

\begin{remember}[$4(x + 3)$ is $4 \times 16$ with the 16 hidden]
  \noindent
  \begin{minipage}[c]{0.40\linewidth}
    \renewcommand{\arraystretch}{1.9}%
    \begin{tabular}{|c|c|c|}
      \hline
      $\times$ & $x$  & $3$ \\ \hline
      $4$      & $4x$ & $12$ \\ \hline
    \end{tabular}%
    \renewcommand{\arraystretch}{1.45}
  \end{minipage}%
  \begin{minipage}[c]{0.58\linewidth}
    $4(x + 3) = \mathbf{4x + 12}$

    \smallskip
    $4 \times x$ is \key{4x} --- you cannot finish it, so you write it and move on.
    $4 \times 3$ is 12. Then \key{add the boxes}, the same as always.
  \end{minipage}
\end{remember}

\vspace{4pt}
\noindent
\begin{minipage}[c]{0.30\textwidth}\gridone{x}{5}{2}\end{minipage}%
\begin{minipage}[c]{0.20\textwidth}\textbf{(a)} $2(x + 5)$ \\[4pt] $= \blank[1.8cm]$\end{minipage}%
\begin{minipage}[c]{0.30\textwidth}\gridone{m}{4}{3}\end{minipage}%
\begin{minipage}[c]{0.20\textwidth}\textbf{(b)} $3(m + 4)$ \\[4pt] $= \blank[1.8cm]$\end{minipage}

\vspace{11pt}
\noindent
\begin{minipage}[c]{0.30\textwidth}\gridone{y}{2}{6}\end{minipage}%
\begin{minipage}[c]{0.20\textwidth}\textbf{(c)} $6(y + 2)$ \\[4pt] $= \blank[1.8cm]$\end{minipage}%
\begin{minipage}[c]{0.30\textwidth}\gridone{k}{7}{5}\end{minipage}%
\begin{minipage}[c]{0.20\textwidth}\textbf{(d)} $5(k + 7)$ \\[4pt] $= \blank[1.8cm]$\end{minipage}

% -- F14 --------------------------------------------------------------------
\qhead[cbgreen]{F14}{Fill the boxes}

\noindent No grid. Write the two numbers straight into the boxes.

\vspace{4pt}
\begin{multicols}{2}
\begin{enumerate}[label=(\alph*), itemsep=13pt]
  \item $3(a + 2) = \nbox a + \nbox$
  \item $5(b + 4) = \nbox b + \nbox$
  \item $7(c + 1) = \nbox c + \nbox$
  \item $2(d + 9) = \nbox d + \nbox$
  \item $4(e + 6) = \nbox e + \nbox$
  \item $8(f + 3) = \nbox f + \nbox$
\end{enumerate}
\end{multicols}

% -- F15 --------------------------------------------------------------------
\qhead[cbgreen]{F15}{A minus inside}

\noindent The minus sign belongs to the number after it, and it goes into the brackets
with that number.

\begin{remember}[Same grid, one sign]
  \hspace*{1.2em}$3(x - 2) \;=\; 3x - 6$ \qquad
  \hspace*{1.2em}$5(y - 4) \;=\; 5y - 20$
\end{remember}

\vspace{2pt}
\begin{multicols}{2}
\begin{enumerate}[label=(\alph*), itemsep=11pt]
  \item $4(x - 3) = \blank[2.4cm]$
  \item $6(m - 2) = \blank[2.4cm]$
  \item $2(k - 7) = \blank[2.4cm]$
  \item $9(n - 1) = \blank[2.4cm]$
  \item $3(p - 5) = \blank[2.4cm]$
  \item $7(q - 2) = \blank[2.4cm]$
\end{enumerate}
\end{multicols}

% -- F16 --------------------------------------------------------------------
% F10 again, with a bracket around it. Same division, same answer.
\qhead[cbgreen]{F16}{Find the number outside}

\noindent This is \textbf{F10} again, with a bracket around it. Divide, exactly as you
did there.

\vspace{4pt}
\begin{multicols}{2}
\begin{enumerate}[label=(\alph*), itemsep=13pt]
  \item $\nbox (x + 2) = 5x + 10$
  \item $\nbox (y + 3) = 4y + 12$
  \item $\nbox (m + 1) = 7m + 7$
  \item $\nbox (k + 5) = 6k + 30$
  \item $\nbox (p + 4) = 3p + 12$
  \item $\nbox (n + 6) = 8n + 48$
\end{enumerate}
\end{multicols}

% -- F17 --------------------------------------------------------------------
% The last question, and the sheet in miniature: the arithmetic from Part 1
% proves the algebra from Part 4, in the same grid.
\qhead[cbgreen]{F17}{Check it with a number}

\begin{yourturn}[This one proves the whole sheet]
  You wrote \; $4(x + 3) = 4x + 12$ \; in F13. Now \key{put $x = 10$ into both
  sides} and see whether they agree.
\end{yourturn}

\vspace{2pt}
\begin{enumerate}[label=(\alph*), itemsep=11pt]
  \item The \textbf{left} side. \quad $x + 3$ is $10 + 3 = \nbox$, \; so
        $4(x+3)$ is $4 \times \nbox = \nbox$
  \item The \textbf{right} side. \quad $4x$ is $4 \times 10 = \nbox$, \; so
        $4x + 12$ is $\nbox + 12 = \nbox$
  \item Did the two sides agree? \hfill \blank[3.0cm]
  \item You have just worked out $4 \times 13$ twice, in two different ways. Which
        two numbers did you add together the second time?
        \hfill \blank[4.0cm]
\end{enumerate}

\vspace{7pt}
\noindent \textbf{(e)} Now do the whole thing again on your own, with
\; \expr{3(m + 4)} \; and \; $m = 20$.

\vspace{4pt}
\begin{enumerate}[label=(\roman*), itemsep=9pt]
  \item Expand it. \hfill \blank[4.4cm]
  \item The \textbf{left} side: work out $m + 4$ first, then multiply by 3.
        \hfill \blank[4.4cm]
  \item The \textbf{right} side: put $m = 20$ into your expanded answer.
        \hfill \blank[4.4cm]
  \item Did they agree? \hfill \blank[4.4cm]
\end{enumerate}

% ===========================================================================
% ANSWER KEY -- TEACHER COPY ONLY
% ===========================================================================
\ifkey
\clearpage

\hwsection[cbred]{Answer Key --- Teacher Copy}

\linespread{1.0}\selectfont
\setlist[enumerate]{itemsep=2pt, topsep=3pt, leftmargin=*}
\setlist[itemize]{itemsep=2pt, topsep=3pt, leftmargin=*}

\qhead[cbred]{F1}{Tens and ones}
\noindent (b) $20+3$ \quad (c) $60+5$ \quad (d) $10+8$ \quad (e) $50+2$ \quad
(f) $90+6$ \quad (g) $30+4$ \quad (h) $70+0$ \quad (i) $80+9$

\smallskip
\noindent{\small (h) is the row to check. $70 = 70 + 0$, and a student who cannot write
it has been splitting by copying the digits rather than by reading the place value.
Accept ``70 and nothing''; do not accept 7 and 0.}

\qhead[cbred]{F2}{Times a ten}
\noindent (a) 120 \quad (b) 300 \quad (c) 280 \quad (d) 540 \quad (e) 160 \quad
(f) 270 \\
(g) 600 \quad (h) 2000 \quad (i) 4200 \quad (j) 2400 \quad (k) 2500 \quad (l) 1800

\smallskip
\noindent{\small (g) to (l) are the ones that decide whether F5 works, because the
top-left box of every two-by-two grid is a ten times a ten and is \emph{much} the
biggest number in the grid. A student who writes 60 for (g) will lose a whole place
value on every question in Part 2 and their answers will be about ten times too small.
That symptom is worth memorising: \textbf{answers ten times too small means F2, not
F5.}}

\qhead[cbred]{F3}{One row: two digits by one digit}
\noindent (a) $140 + 21 = \mathbf{161}$ \quad (b) $240 + 48 = \mathbf{288}$ \\
(c) $400 + 32 = \mathbf{432}$ \quad (d) $270 + 63 = \mathbf{333}$

\smallskip
\noindent{\small Mark the boxes, not the total. The point of a printed grid is that
every step is visible, so a wrong answer here can always be traced to one box. \\
Both boxes filled and the addition wrong is an arithmetic slip. One box filled is the
error this sheet exists to fix.}

\qhead[cbred]{F4}{Now without the printed grid}
\noindent 204 \quad 576 \quad 595 \quad 414 \\
340 \quad 372 \quad 456 \quad 448 \\
702 \quad 570 \quad 343 \quad 688

\smallskip
\noindent{\small In order: $6 \times 34$, $8 \times 72$, $7 \times 85$, $9 \times 46$;
$5 \times 68$, $4 \times 93$, $8 \times 57$, $7 \times 64$; $9 \times 78$,
$6 \times 95$, $7 \times 49$, $8 \times 86$. \\
The instruction is to \emph{draw the grid}, and it is worth insisting on for another
week even from students who no longer need it --- Part 4 asks for the same picture with
a letter in it, and a student who has stopped drawing it has nothing to fall back on
when the box will not multiply.}

\qhead[cbred]{F5}{Four boxes instead of two}
\noindent (a) $300+20+150+10 = \mathbf{480}$ \quad
(b) $800+120+120+18 = \mathbf{1058}$ \\
(c) $1500+210+200+28 = \mathbf{1938}$ \quad
(d) $2400+320+120+16 = \mathbf{2856}$ \\
(e) $1400+420+80+24 = \mathbf{1924}$ \quad
(f) $2400+560+150+35 = \mathbf{3145}$

\smallskip
\noindent{\small Count the filled boxes before reading the answer. Four numbers means
the method was followed; three means one box was missed, which is the single commonest
error in two-by-two multiplication and is invisible in the total. \\
In (b), $23$ splits as $20+3$ and both rows have to reach both columns. The box that
gets forgotten is nearly always the bottom-right one --- the small ones times the small
ones --- because it looks least important and is the easiest to do.}

\qhead[cbred]{F6}{Draw the grid yourself}
\noindent 312 \quad 735 \quad 752 \\
1768 \quad 1701 \quad 3510 \\
4558 \quad 3572 \quad 4800 \\
1073 \quad 2688 \quad 2158

\smallskip
\noindent{\small In order: $24 \times 13$, $35 \times 21$, $47 \times 16$;
$52 \times 34$, $63 \times 27$, $78 \times 45$; $86 \times 53$, $94 \times 38$,
$75 \times 64$; $29 \times 37$, $56 \times 48$, $83 \times 26$. \\
$75 \times 64 = 4800$ is a good one to put on the board at the end, because the grid
gives $4200 + 300 + 280 + 20$ and the round answer looks like a coincidence until the
four boxes are shown.}

\qhead[cbred]{F7}{Use it}
\noindent 1. \; $24 \times 36 = \mathbf{864}$ pencils. \quad
2. \; $18 \times 15 = \mathbf{270}$ chairs. \\
3. \; $47 \times 23 = \mathbf{1081}$ seats. \quad
4. \; $18 \times 25 = \mathbf{450}$ rubber ducks.

\smallskip
\noindent{\small None of the four says \emph{multiply}. The words doing the work are
\emph{every box holds}, \emph{in each row}, \emph{in total} and \emph{every tank holds},
which are all in the word help. A student who adds $24+36$ has read the numbers and not
the sentence --- that is a reading error, and it should be named as one out loud rather
than marked as arithmetic. \\
3 is the hardest arithmetic on the sheet and the grid earns its keep:
$800+140+210+21$. \\
4 is meant to be read completely straight. If the class laughs, that is fine; the
question is still $18 \times 25$, and $450$ rubber ducks is the correct answer.}

\qhead[cbred]{F8}{The same three numbers, four ways}
\noindent Row 2: $7 \times 8 = 56$, $8 \times 7 = 56$, $56 \div 7 = 8$,
$56 \div 8 = 7$. \\
Row 3: $8 \times 9 = 72$, $9 \times 8 = 72$, $72 \div 8 = \mathbf{9}$,
$72 \div 9 = 8$. \\
Row 4: $8 \times 12 = 96$, $12 \times 8 = 96$, $96 \div 8 = 12$,
$96 \div 12 = \mathbf{8}$.

\smallskip
\noindent{\small Rows 3 and 4 are filled in \emph{backwards}, from a division to the
multiplication that made it, which is the whole idea Part 3 is for: a division question
is a multiplication question you already know the answer to. \\
This is F10 and F16 in advance. If a student can do this table, neither of those is new
work.}

\qhead[cbred]{F9}{Divide}
\noindent (a) 12 \quad (b) 12 \quad (c) 12 \quad (d) 15 \quad (e) 32 \quad (f) 23
\quad (g) 24 \quad (h) 36 \\
(i) 8 r 2 \quad (j) 16 r 4 \quad (k) 9 r 4 \quad (l) 35 r 3

\smallskip
\noindent{\small (a), (b) and (c) are all 12 on purpose. Say so afterwards --- a student
who trusted their arithmetic and got three identical answers will have assumed they were
wrong, and that assumption is worth naming. \\
(f) $161 \div 7$ is the hardest of the exact ones. Counting up in sevens is a perfectly
good method and should not be marked down.}

\qhead[cbred]{F10}{Find the missing number}
\noindent (a) 9 \quad (b) 12 \quad (c) 12 \quad (d) 12 \quad (e) 5 \quad (f) 23

\smallskip
\noindent{\small Every one of these is found by dividing --- $54 \div 6$, $84 \div 7$,
$108 \div 9$, $96 \div 8$, $60 \div 12$, $115 \div 5$. Say that sentence out loud when
you hand the sheet back, because F16 is this question with a bracket drawn round it and
the connection is the point. \\
A student who gets these by guessing and checking has still got them right. Ask for the
division as well, once.}

\qhead[cbred]{F11}{Share it out}
\noindent 1. \; $144 \div 12 = \mathbf{12}$ pencils each.

\smallskip
\noindent 2. \; $90 \div 16 = 5$ r 10, so \textbf{6 buses} are needed.

\smallskip
\noindent 3. \; $216 \div 8 = \mathbf{27}$ stickers each. \quad
4. \; $175 \div 18 = \mathbf{9}$ each, with \textbf{13 left over}.

\smallskip
\noindent{\small \textbf{Question 2 is the one worth discussing.} The arithmetic gives 5
remainder 10 and the answer is 6, because ten students cannot be left in the car park.
Expect 5, and expect ``5.6''. \\
Then put 2 and 4 side by side. Both divisions leave a remainder and the two remainders
are treated in opposite ways: the ten students need a sixth bus, and the thirteen spare
ducks simply stay in the box. Nothing in the arithmetic tells you which --- only the
sentence does. That is the whole reason these are word problems and not sums. \\
4 is meant to be read straight, and 9 r 13 is the full answer; a student who writes 9 has
answered only half the question.}

\qhead[cbred]{F12}{A letter is a number you do not know yet}
\noindent (b) $9b$ \quad (c) $5c$ \quad (d) $4x$ \quad (e) $2m$ \quad (f) $5y$ \quad
(g) $2p+3q$ \textbf{finished} \quad (h) $5k$ \quad (i) $4n+7$ \textbf{finished}

\smallskip
\noindent{\small (d), (e) and (f) are the invisible 1, three times. (d) comes out $5$ if
the $x$ is read as nothing; (f) comes out $6$ the same way. \\
(g) and (i) are the two to leave alone, and they are the reason the letter work comes
before the brackets: a student who turns $4n+7$ into $11n$ will turn $4(x+3)$ into
$4x+12$ and then into $16x$.}

\qhead[cbred]{F13}{One box has a letter in it}
\noindent (a) $2x+10$ \quad (b) $3m+12$ \quad (c) $6y+12$ \quad (d) $5k+35$

\smallskip
\noindent{\small The left box is the whole lesson. $2 \times x$ cannot be worked out, so
it is written as $2x$ and left --- and a student who stalls there has not misunderstood
multiplication, they have been taught that a box must contain a number. Say plainly that
this box does not. \\
Compare (c) with F3(d) if anyone doubts the grids are the same picture: identical
layout, identical instruction, one box that happens to be a letter.}

\qhead[cbred]{F14}{Fill the boxes}
\noindent (a) $3a+6$ \quad (b) $5b+20$ \quad (c) $7c+7$ \quad (d) $2d+18$ \quad
(e) $4e+24$ \quad (f) $8f+24$

\smallskip
\noindent{\small (c) is the one that looks wrong and is not: $7 \times 1 = 7$, so both
boxes hold a 7. Expect it to be left blank. \\
The printed boxes are doing real work here. A blank square in front of the letter is
much harder to skip than an empty line, and skipping the second number ---
$3(a+2) = 3a+2$ --- is the headline error of Maths 2.4.}

\qhead[cbred]{F15}{A minus inside}
\noindent (a) $4x-12$ \quad (b) $6m-12$ \quad (c) $2k-14$ \quad (d) $9n-9$ \quad
(e) $3p-15$ \quad (f) $7q-14$

\smallskip
\noindent{\small (d) is the invisible 1 arriving inside a bracket: $9 \times 1 = 9$, so
the answer is $9n-9$ and not $9n-1$. It is F14(c) with a minus sign, and a student who
got that one right and this one wrong was copying rather than multiplying.}

\qhead[cbred]{F16}{Find the number outside}
\noindent (a) 5 \quad (b) 4 \quad (c) 7 \quad (d) 6 \quad (e) 3 \quad (f) 8

\smallskip
\noindent{\small Found by dividing: $10 \div 2$, $12 \div 3$, $7 \div 1$, $30 \div 5$,
$12 \div 4$, $48 \div 6$.
Or by dividing the other term --- $5x \div x$ --- which is also right and is worth
hearing if a student offers it. \\
Point back at F10 before they start. The two questions are the same question and the
sheet says so; a student who does F10 confidently and stalls here has been stopped by
the bracket, not by the arithmetic.}

\qhead[cbred]{F17}{Check it with a number}
\noindent (a) $10+3 = \mathbf{13}$, \; so $4 \times \mathbf{13} = \mathbf{52}$. \\
(b) $4 \times 10 = \mathbf{40}$, \; so $\mathbf{40} + 12 = \mathbf{52}$. \\
(c) \textbf{Yes} --- both sides are 52. \\
(d) \textbf{40 and 12.}

\smallskip
\noindent (e) (i) $3m+12$ \; (ii) $m+4 = 24$, and $3 \times 24 = \mathbf{72}$ \;
(iii) $3 \times 20 = 60$, and $60 + 12 = \mathbf{72}$ \; (iv) \textbf{yes}.

\smallskip
\noindent{\small (e) is (a) to (d) with the scaffolding taken away, and it is the one
that shows whether the method transferred. The two boxes this time are \textbf{60 and
12}, which is the grid for $3 \times 24$. \\
\noindent{\small \textbf{This is the sheet in one question.} The left side is
$4 \times 13$ done as a single multiplication. The right side is $4 \times 13$ done as
$40 + 12$ --- which is the grid from F3, with the same two boxes, arrived at from the
algebra. \\
(d) is the part to ask for out loud. A student who says ``40 and 12'' and then sees that
those are the two boxes of $4 \times 13$ has understood why expanding a bracket is not a
new rule but the multiplication they have done all year, written down before the letter
is known. \\
If there is time for one thing from this sheet at the board, do this question and draw
the grid for $4 \times 13$ beside it.}

\fi

\end{document}
